% ======================================================================
%  udt — plik startowy: STRESZCZENIE WYKŁADU (tryb report)
%
%  Tryb report nadaje się nie tylko do sprawozdań: z \doctype staje się
%  dwukolumnowym streszczeniem albo notatką. Ten plik pokazuje to, czego
%  potrzeba w notatkach z wykładu: środowiska twierdzeń, dowody, uwagi
%  na marginesie i matematykę (pochodne, wektory, bra-ket, macierze).
%  Pola nagłówka, rysunki, tabele i bibliografia: zob. sprawozdanie.tex.
%  Warianty i opcja chem: zob. skrypt.tex.
%
%  Kompilacja:  pdflatex → pdflatex
%  (drugi przebieg: odsyłacze i pionowe kreski dowodów)
% ======================================================================
\documentclass[a4paper, polish, 11pt]{scrartcl}
\usepackage{udt-report}

% ----------------------------------------------------------------------
%  NAGŁÓWEK
%  \doctype zmienia rodzaj dokumentu; \exercise dopisuje do niego numer
%  (dowolny tekst). Tu: „Streszczenie wykładu 1: część I”, nad tytułem
%  i w paginie. Bez \exercise zostaje sam \doctype, np. „Notatki”.
%  \date{} usuwa wiersz „Oddano: …”, zbędny w notatkach.
% ----------------------------------------------------------------------
\doctype{Streszczenie wykładu}
\exercise{1: część I}
\title{Tytuł wykładu}
\author{Imię Nazwisko}
\date{}

\begin{document}
\maketitle

\begincolumns

\section{Środowiska twierdzeń}\label{sec:srodowiska}

% ----------------------------------------------------------------------
%  ŚRODOWISKA TWIERDZEŃ
%  Każde jest kolorowym blokiem z paskiem po lewej. Wspólny licznik
%  (w trybie report 1, 2, … w całym dokumencie), więc „Definicja 3”
%  wskazuje dokładnie jeden obiekt. Nazwa własna w nawiasie kwadratowym;
%  przy kilku nazwach — wszystkie, najpopularniejsza pierwsza.
%    definition — definicja (prosto, niebieski),
%    theorem    — twierdzenie (kursywa, czerwony),
%    law        — prawo fizyczne (prosto, fioletowy),
%    lemma, fact, corollary — lemat, fakt, wniosek (kursywa, zielony),
%    example    — przykład (prosto, szary),
%    remark     — uwaga (prosto, pomarańczowy).
%  Każde ma bliźniaka z gwiazdką (definition*, …) — ten sam wygląd, bez
%  numeru. Odsyłacz: \cref{def:hermit} → „definicja 1”.
%  Opcja plain wyłącza kolorowe bloki (klasyczny wygląd amsthm).
% ----------------------------------------------------------------------
\begin{definition}[Macierz hermitowska]\label{def:hermit}
Macierz kwadratową $A$ nazywa się hermitowską, jeżeli $A^\dagger = A$, czyli $A_{jk} = A_{kj}^*$ dla wszystkich $j, k$.
\end{definition}

\begin{theorem}[Twierdzenie spektralne]\label{tw:spektralne}
Każda macierz hermitowska $A$ ma rzeczywiste wartości własne i daje się zapisać jako $A = Q\Lambda Q^\dagger$, gdzie $Q$ jest unitarna, a $\Lambda$ diagonalna.
\end{theorem}

% proof — dowód z cienką szarą kreską po lewej i znakiem końca ∎.
% Nazwa w osobnym wierszu; inna nazwa: \begin{proof}[Szkic dowodu].
% Kreska pojawia się po drugim przebiegu kompilacji.
\begin{proof}
Dla wektora własnego $\ket{v}$ z wartością $\lambda$ mamy $\braket*[3]{v}{A}{v} = \lambda\braket{v}{v}$. Lewa strona jest rzeczywista, bo $A^\dagger = A$, a $\braket{v}{v} > 0$, więc $\lambda \in \mathbb{R}$.
\end{proof}

\begin{example}
Macierz $\sigma_y = \begin{pmatrix} 0 & -i \\ i & 0 \end{pmatrix}$ jest hermitowska (\cref{def:hermit}).
\end{example}

\begin{remark}
\Cref{tw:spektralne} pozwala sprowadzić każde równanie Schrödingera do układu niezależnych oscylatorów.
\end{remark}

% aside — „na marginesie”: informacja poboczna (historia, ciekawostka,
% zastosowanie), którą można pominąć bez utraty wątku. Bez numeru,
% mniejsze pismo, szary blok. Tytuł opcjonalny.
\begin{aside}[Skąd nazwa?]
Charles Hermite badał takie formy kwadratowe w XIX wieku, zanim pojawiła się mechanika kwantowa.
\end{aside}

\section{Matematyka}\label{sec:matematyka}

% ----------------------------------------------------------------------
%  POCHODNE (pakiet derivative) i RÓŻNICZKI (fixdif)
%    \odv{f}{t}        pochodna zwyczajna,  \odv[2]{x}{t}  druga,
%    \odv{}{t}         sam operator d/dt,
%    \pdv{f}{x}        cząstkowa,           \pdv{f}{x,y}   mieszana,
%    \dv               alias \odv (nawyk z pakietu physics),
%    \d x, \dd x       prosta różniczka z poprawnym odstępem.
%  Pochodne po czasie kropką: \dot{x}_k, \ddot{q}_2.
% ----------------------------------------------------------------------
Równania Hamiltona:
\begin{equation}
	\dot{x}_k = \pdv{\bar H}{p_k}, \qquad \dot{p}_k = -\pdv{\bar H}{x_k}
	\label{eq:hamilton}
\end{equation}
a praca siły na drodze $\Gamma$ to $W = \int_\Gamma \va{F} \vdot \d\va{r}$.

% ----------------------------------------------------------------------
%  WEKTORY
%    \vb{r}  pogrubiony prosty,    \vb*{r}  pogrubiony pochyły,
%    \va{F}  ze strzałką (prosty), \va*{F}  ze strzałką (pochyły),
%    \vu{n}  wersor (daszek),      \vu*{n}  wersor pochyły,
%    \vdot, \cross — iloczyn skalarny i wektorowy,
%    \grad, \div, \curl, \laplacian — operatory nabla
%    (symbol dzielenia ÷ to \divsymbol).
% ----------------------------------------------------------------------
Moment siły $\va{M} = \vb{r} \cross \va{F}$, rzut na oś $\vu{n}$: $M_n = \vu{n} \vdot \va{M}$. Pole zachowawcze: $\va{F} = -\grad V$, $\curl\va{F} = 0$.

% ----------------------------------------------------------------------
%  NAWIASY (physics2)
%    \ab(...), \ab[...], \ab\{...\} — nawiasy dopasowane do zawartości
%                                     (zamiast \qty(...) z physics),
%    \abs{x}, \norm{v}               — moduł i norma,
%    \eval{f}_{a}^{b}                — wartość w granicach.
%  Uwaga: \qty to w udt polecenie siunitx (liczba z jednostką),
%  np. \qty{9.81}{\metre\per\second\squared}.
% ----------------------------------------------------------------------
Energia: $\bar H = \sum_{k} \ab(\frac{p_k^2}{2m_k} + \frac{m_k\omega_k^2 x_k^2}{2})$, a $\abs{z_k}^2 = z_k z_k^*$.

% ----------------------------------------------------------------------
%  BRA-KET (physics2, składnia z klamrami)
%    \ket{\psi}, \bra{\psi}          — ket, bra,
%    \braket{\phi}{\psi}             — iloczyn skalarny ⟨φ|ψ⟩,
%    \braket[3]{\phi}{A}{\psi}       — element macierzowy ⟨φ|A|ψ⟩,
%    \ketbra{\psi}{\phi}             — |ψ⟩⟨φ|.
%  Nawiasy i kreski same się skalują. Gdy wysoki element (kropka,
%  krzyżyk †) powiększa je ponad potrzebę, wersja z gwiazdką ich nie
%  skaluje: \ket*{\dot\psi}, \braket*[3]{v}{A^\dagger A}{w}.
% ----------------------------------------------------------------------
Równanie Schrödingera $i\hbar\,\odv{}{t}\ket{\psi} = H\ket{\psi}$ zachowuje normę $\braket{\psi}{\psi}$, a energia to $\bar H = \braket[3]{\psi}{H}{\psi}$. Pochodna: $\ket*{\dot\psi} = -\frac{i}{\hbar}H\ket{\psi}$, rzutowanie: $P = \ketbra{\psi}{\psi}$.

% ----------------------------------------------------------------------
%  MACIERZE
%    pmatrix, bmatrix, vmatrix (amsmath) — nawiasy ( ), [ ], | |,
%    \pdiagmat{a, b, c}                — macierz diagonalna (zera poza
%                                        przekątną; \diagmat bez nawiasów),
%    \pxmat{a}{3}{3}                   — macierz ogólna a_{jk},
%    \mathds{1}                        — macierz jednostkowa,
%    \Tr, \sgn                         — ślad, znak.
%  Szeroki wzór w kolumnie: \setlength{\arraycolsep}{2pt} wewnątrz
%  equation zwęża odstępy między kolumnami macierzy.
% ----------------------------------------------------------------------
Macierz częstości i ogólna macierz $3 \times 3$:
\begin{equation}
	\setlength{\arraycolsep}{2pt}
	\Omega = \pdiagmat{\omega_1, \omega_2, \omega_3}, \qquad A = \pxmat{a}{3}{3}
	\label{eq:macierz}
\end{equation}
przy czym $\Tr\Omega = \omega_1 + \omega_2 + \omega_3$ i $A^0 = \mathds{1}$.

% ----------------------------------------------------------------------
%  WZORY WIELOWIERSZOWE
%  W wąskiej kolumnie długi wzór łamie się przez aligned wewnątrz
%  equation (jeden numer dla całości). Przy łamaniu na znaku równości
%  znak powtarza się w nowym wierszu.
%  Wzory wystawione nie mają kropki ani przecinka na końcu.
% ----------------------------------------------------------------------
\begin{equation}
	\begin{aligned}
		e^X &= \sum_{m=0}^{\infty} \frac{X^m}{m!} \\
		&= \mathds{1} + X + \frac{X^2}{2!} + \ldots
	\end{aligned}
	\label{eq:eksponenta}
\end{equation}

% Układ równań z klamrą:
\begin{equation}
	\left\{
	\begin{aligned}
		\dot{x}_k &= v_k \\
		\dot{v}_k &= -\omega_k^2\,x_k
	\end{aligned}
	\right.
	\label{eq:uklad}
\end{equation}

% Podpis pod wyrazem: \underbrace{\hbar\Omega}_{H}.
Zatem $i\hbar\,\odv{}{t}\ket{z} = \underbrace{\hbar\Omega}_{H}\ket{z}$ (\cref{eq:uklad,eq:eksponenta}).

\end{document}
